\documentclass[11pt,a4paper]{article}

\usepackage[margin=30mm]{geometry}
\usepackage[T1]{fontenc}
\usepackage{lmodern}
\usepackage{microtype}
\usepackage{amsmath,amssymb,amsthm,mathtools}
\usepackage{booktabs}
\usepackage{enumitem}
\usepackage[colorlinks=true,linkcolor=blue,citecolor=blue,urlcolor=blue]{hyperref}
\usepackage[nameinlink,capitalise]{cleveref}

\newcommand{\kk}{\Bbbk}
\newcommand{\hK}{\mathbb K}
\newcommand{\ot}{\otimes}
\newcommand{\hot}{\widehat\otimes}
\newcommand{\tri}{\mathbin{\triangleleft}}
\newcommand{\id}{\operatorname{id}}
\newcommand{\ev}{\operatorname{ev}}
\newcommand{\coev}{\operatorname{coev}}
\newcommand{\Writhe}{\operatorname{wr}}
\newcommand{\End}{\operatorname{End}}
\newcommand{\ao}{\mathrm{ao}}
\newcommand{\rc}{\mathrm{rc}}
\newcommand{\cen}{\mathrm{cen}}
\newcommand{\rootpart}{\mathrm{root}}

\newtheorem{theorem}{Theorem}[section]
\newtheorem{proposition}[theorem]{Proposition}
\newtheorem{lemma}[theorem]{Lemma}
\newtheorem{corollary}[theorem]{Corollary}
\theoremstyle{definition}
\newtheorem{definition}[theorem]{Definition}
\theoremstyle{remark}
\newtheorem{remark}[theorem]{Remark}

\title{The \(V_n\) Invariants as Colored Links--Gould Invariants:\\
A Root--Center Approach}
\author{Jiuhe Liu}
\date{}

\begin{document}

\maketitle

\begin{abstract}
The knot invariants \(V_n\) arise from a rank-two Nichols algebra and a
\(4n\)-dimensional right Yetter--Drinfeld module, whereas the colored
Links--Gould invariants are defined from typical
\(U_q(\mathfrak{sl}(2|1))\)-modules.  We prove that these two constructions
agree, up to the mirror and parameter inversion forced by the right-module
convention.  The proof is structural.  We first realize the braid operator
\(T_n\) as a gauge transform of a canonical super
Yetter--Drinfeld braiding.  We then construct the relevant right--right
paired double, prove that its Hopf pairing is perfect in every root degree,
and obtain its completed universal \(R\)-matrix.  An explicit Abelian
Drinfeld twist separates this double into an all-odd
\(U_\hbar(\mathfrak{sl}(2|1))\) root factor and a commutative central
factor.  Under this factorization, the Nichols module becomes a typical
all-odd highest-weight module tensored with a one-dimensional central
module.  Finally, we transport duality, all four oriented crossings,
partial transposes, and writhe normalization.  For every oriented knot
\(\mathcal K\), every \(n\geq 1\), and every
\(\beta\neq 0,-1\), the result is
\[
 V_{n,\mathcal K}\!\left(Q^{2n\beta+n},Q^2\right)
 =
 LG_{\mathcal K}^{(n)}\!\left(Q^{-n\beta},Q^{-1}\right)
 =
 LG_{\overline{\mathcal K}}^{(n)}\!\left(Q^{n\beta},Q\right).
\]
In particular, the scalar-identity property conjectured for the
endomorphism-valued \(V_n\) construction follows from the simplicity of
the corresponding typical module.
\end{abstract}

\tableofcontents

\section{Introduction}

The invariants \(V_n\) of Garoufalidis and Kashaev are extracted from a
family of braid-form Yang--Baxter operators \(T_n\) acting on
\(4n\)-dimensional spaces \(Y_n\).  The construction is rooted in a
rank-two Nichols algebra and is strikingly explicit.  The colored
Links--Gould invariant, by contrast, is built from typical modules over a
quantum deformation of \(\mathfrak{sl}(2|1)\).  The two theories therefore
begin with rather different algebraic input.  The purpose of this paper is
to explain why they nevertheless produce the same knot invariant.

The connection is not obtained by a direct comparison of two large
matrices.  Instead, we expose a sequence of structural identifications:
\[
\begin{gathered}
\text{Nichols braid operator }T_n
\longleftrightarrow
\text{super Yetter--Drinfeld braiding}\\
\longleftrightarrow
\text{braiding of a completed paired double}\\
\xrightarrow{\ \text{root--center twist}\ }
U_\hbar^{\Pi_{\ao}}(\mathfrak{sl}(2|1))
\hot \mathcal Z_{\cen}.
\end{gathered}
\]
The \(4n\)-dimensional object splits at the last stage as
\[
 \mathcal Y_n\cong W_{n,\beta}^{\ao}\boxtimes L_{n,\beta},
\]
where \(W_{n,\beta}^{\ao}\) is a simple all-odd highest-weight module and
\(L_{n,\beta}\) is one-dimensional and central.  The central line
contributes a computable crossing scalar.  It disappears on balanced
normal long-knot diagrams and, more generally, under the standard writhe
normalization.

There are two points at which a braid-level comparison could fail to
produce a knot invariant.  First, the Hopf pairing used to form the double
must be well defined, nondegenerate in the relevant completion, and strong
enough to produce a genuine universal \(R\)-matrix.  We prove this using
explicit alternating-word bases and root-degree pairing matrices.  Second, a
conjugacy on \(Y_n\otimes Y_n\) does not by itself control cups, caps,
mixed orientations, partial transposes, or curls.  We therefore transport
the entire rigid tangle calculus and compare the transported duality maps
with the ordinary finite-dimensional ones.

\subsection{Related work and scope}

The original Links--Gould invariant was constructed from the
four-dimensional family of typical modules over
\(U_q(\mathfrak{gl}(2|1))\), using a tensor-state model for
\((1,1)\)-tangles \cite{DeWitKauffmanLinks}.  This is an instance of the
Reshetikhin--Turaev construction from a ribbon category
\cite{ReshetikhinTuraev}.  Because typical modules for Lie superalgebras
have vanishing ordinary quantum dimension, the natural closed-link theory
is more conveniently expressed using a renormalized invariant or modified
trace.  For \(\mathfrak{sl}(2|1)\), this point of view was developed by
Geer and Patureau--Mirand \cite{GeerPatureau}, and it underlies the modern
colored Links--Gould family.

The relevant algebraic background has several distinct layers.  The
classification and root theory of basic classical Lie superalgebras go
back to Kac \cite{Kac}, while changes of Borel through isotropic odd roots
belong to the generalized-root-system framework of Serganova
\cite{Serganova}.  Quantized enveloping superalgebras and their universal
\(R\)-matrices were constructed explicitly by Yamane \cite{Yamane}.
Etingof--Kazhdan quantization in the super setting and its relation to
Drinfeld--Jimbo presentations, including arbitrary Cartan matrices, were
studied by Geer \cite{GeerEK,GeerCartan}.  Geer's comparison with the
Kontsevich integral identifies the corresponding quantum-group and
weight-system invariants for Lie superalgebras of type A--G
\cite{GeerKontsevich}.  These general results supply essential structural
background, but they do not by themselves identify the particular
right--right completed double, module gauges, duality conventions, and
writhe normalization used below; those interfaces are treated explicitly
in this paper.

On the Nichols-algebra side, the spaces \(Y_n\) and their braid operators
were introduced by Garoufalidis and Kashaev in their construction of
multivariable knot polynomials from braided Hopf algebras with
automorphisms \cite{GK}.  Garoufalidis, Harper, Kashaev, Kohli, and Wagner
subsequently defined the colored Links--Gould sequence from
\(4n\)-dimensional typical modules, established its cabling and modified
trace framework, and formulated the equality with the \(V_n\)-sequence as
a conjecture; in that work the comparison was proved for \(n=2\)
\cite[Conjecture~1.4 and Theorem~1.5]{ColoredLG}.  Thus the contribution
claimed here is not the definition of either family.  It is the proposed
all-\(n\) structural bridge: the realization of the Nichols operator in a
completed paired double, its root--center separation, the identification
of the coloring object, and the transport of the full oriented rigid
tangle data needed to compare the normalized knot scalars.

We work with a formal quantum parameter, away from the atypical values
\(\beta=0,-1\).  The comparison concerns normalized oriented-knot
invariants, computed from normal \((1,1)\)-tangles; it does not assert an
equivalence of the entire module categories.  The completion conventions
are specified in \cref{sec:conventions,sec:completion}.

\subsection{Main results}

Let \(\kk\) be a field of characteristic zero,
\[
 R=\kk[[\hbar]],\qquad \hK=\kk((\hbar)),\qquad Q=e^\hbar.
\]
The invariant \(LG_{\mathcal K}^{(n)}(q^\alpha,q)\) is normalized using the
\(4n\)-dimensional typical distinguished-Borel module
\(V_0(n-1,\alpha)\), with value \(1\) on the unknot.

\begin{theorem}[Main theorem]\label{thm:main}
For every \(n\geq 1\), every \(\beta\in\kk\setminus\{0,-1\}\), and every
oriented knot \(\mathcal K\),
\begin{equation}\label{eq:main}
 V_{n,\mathcal K}\!\left(Q^{2n\beta+n},Q^2\right)
 =
 LG_{\mathcal K}^{(n)}\!\left(Q^{-n\beta},Q^{-1}\right)
 =
 LG_{\overline{\mathcal K}}^{(n)}\!\left(Q^{n\beta},Q\right).
\end{equation}
Moreover, the endomorphism-valued long-knot invariant defining \(V_n\)
is a scalar multiple of the identity on \(\mathcal Y_n\).
\end{theorem}

The algebraic mechanism behind \cref{thm:main} is summarized by the next
two statements.

The proof has two algebraic ingredients.  The completed paired double
has a root-degree perfect Hopf pairing and a universal \(R\)-matrix
(\cref{thm:universal-R}); an Abelian twist separates its root and central
Hopf factors (\cref{thm:root-center,prop:R-factor}).  Under this
separation, the coloring object has the module decomposition above
(\cref{thm:object}).  The braid comparison extends to oriented crossings
and duality, yielding the long-knot transport formula of \cref{thm:rigid}.

\subsection{Outline of the proof}

Section~2 constructs the positive Hopf superalgebra and its finite-dimensional
super Yetter--Drinfeld module.  Section~3 constructs the completed paired
double and its universal \(R\)-matrix.  Section~4 separates the root and
central factors, and Section~5 identifies the coloring module, including
its distinguished-Borel highest weight.  Section~6 compares all oriented
crossings and duality maps.  Section~7 assembles these local identities
into long-knot transport and proves the comparison with the standard
colored Links--Gould invariant.  Appendix~A checks the original
Garoufalidis--Kashaev conventions and records the parameter dictionary;
Appendix~B gives the first root-degree terms of the canonical tensor.

\section{The Nichols module and its braiding}
\subsection{Conventions}\label{sec:conventions}

Unless otherwise stated, algebras, modules, pairings, and linear maps
are over \(\hK\).  Formal quantum groups and exponentials are first
defined over \(R\), with tensor products of their \(R\)-forms
\(\hbar\)-adically completed, and are then extended to \(\hK\).
In particular, \(\hK\hot_R U_\hbar\) means
\(\hK\ot_R U_\hbar\) with the \(R\)-form \(U_\hbar\) already
complete; no further ideal-adic completion after inverting \(\hbar\)
is intended.  Completed tensor products of the resulting
\(\hK\)-algebras are likewise formed from the indicated completed
\(R\)-forms before extension of scalars.  The field \(\hK\) carries
its Laurent series valuation topology, not the ideal-adic topology
of \((\hbar)\subset\hK\), since \(\hbar\) is invertible.
Infinite root-degree products use the additional directed completion
specified in \cref{sec:completion}.  For an even element \(H\),
\[
 Q^H=e^{\hbar H}.
\]
If \(a,b,c,d\) are homogeneous elements of super vector spaces, then
\begin{equation}\label{eq:koszul}
 (a\ot b)(c\ot d)=(-1)^{|b||c|}ac\ot bd.
\end{equation}

Every module in this paper is a right module.  Thus, if
\(\rho(X)(v)=v\tri X\), then
\begin{equation}\label{eq:right-antihom}
 \rho(XY)=\rho(Y)\rho(X).
\end{equation}
Our universal \(R\)-matrix convention is
\begin{equation}\label{eq:R-convention}
 \mathcal R\Delta(a)=\Delta^{\mathrm{op}}(a)\mathcal R.
\end{equation}
It follows from \eqref{eq:right-antihom} that the canonical braiding on
right modules is
\begin{equation}\label{eq:right-braiding}
 c^{\mathcal R}_{V,W}
 =\tau^{\mathrm s}\rho_{V\ot W}(\mathcal R^{-1}),
\end{equation}
where \(\tau^{\mathrm s}\) is the super flip.  This inverse is responsible
for the mirror in \cref{eq:main}.

\subsection{The reduced positive half}
\label{sec:reduced}

Let \(\mathcal H^+\) be the topological Hopf superalgebra over
\(\hK\) obtained by extending scalars from the \(\hbar\)-adic
completion of the \(R\)-superalgebra generated by even elements
\(H_1,H_2\) and odd elements \(x_1,x_2\), subject to
\begin{equation}\label{eq:positive-relations}
 [H_1,H_2]=0,\qquad
 [H_i,x_j]=a_{ij}x_j,\qquad x_1^2=x_2^2=0,
 \qquad
 A=(a_{ij})=\begin{pmatrix}0&1\\1&0\end{pmatrix}.
\end{equation}
No additional quadratic relation is imposed between \(x_1x_2\) and
\(x_2x_1\).  The Hopf structure is
\begin{equation}\label{eq:positive-Hopf}
\begin{aligned}
 \Delta(H_i)&=H_i\ot1+1\ot H_i,
 &\varepsilon(H_i)&=0,
 &S(H_i)&=-H_i,\\
 \Delta(x_i)&=x_i\ot1+Q^{H_i}\ot x_i,
 &\varepsilon(x_i)&=0,
 &S(x_i)&=-Q^{-H_i}x_i.
\end{aligned}
\end{equation}
\subsection[The space Yn]{The space \(\mathcal Y_n\)}

Let \(B_{\hK}\) be the superalgebra generated by \(x_1,x_2\) with only
\(x_1^2=x_2^2=0\).  A bihomogeneous element of bidegree \((a,b)\) has parity
\(\overline{a+b}\).  For nonzero bihomogeneous elements
\(x,y\in B_{\hK}\) with \(\deg x=(a,b)\) and \(\deg y=(c,d)\), define
\begin{equation}\label{eq:chi}
 \chi(x,y)=(-1)^{(a+b)(c+d)}Q^{ad+bc}.
\end{equation}
The full diagonal braiding is the linear map determined by
\[
 \tau_\chi(x\ot y)=\chi(x,y)y\ot x.
\]
In particular,
\[
 \chi(x,x_1)=(-1)^{a+b}Q^b,\qquad
 \chi(x,x_2)=(-1)^{a+b}Q^a.
\]
We use the same notation on the bihomogeneous elements of
\(\mathcal Y_n\subset B_{\hK}\) defined below.  All subsequent uses of
\(\chi\) and all Sweedler formulas are understood on bihomogeneous
components.
Fix \(n\geq1\) and set
\begin{equation}\label{eq:parameters}
 t=Q^{2n\beta+n}.
\end{equation}
Define
\[
 v_n=(x_1x_2)^n+t(x_2x_1)^n.
\]
For \(\beta\neq0,-1\), let \(\mathcal Y_n\) be the span of
\begin{equation}\label{eq:Yn-basis}
\begin{split}
 \{1,v_n\}
 &\cup\{(x_1x_2)^k,(x_2x_1)^k:1\leq k<n\}\\
 &\cup\{(x_1x_2)^kx_1,(x_2x_1)^kx_2:0\leq k<n\}.
\end{split}
\end{equation}
This displayed set is a basis, so \(\dim_{\hK}\mathcal Y_n=4n\).

The parameter and convention dictionary identifying \(\mathcal Y_n\) with
the right Yetter--Drinfeld module \(Y_n\) of
\cite[Proposition~7.4]{GK} is given in
\cref{app:GK-conventions}.  Under this identification, \(T_n\) denotes
precisely the braid-form \(R\)-matrix defined by
\cite[Theorem~3.6 and Definition~7.5]{GK}, after extension of scalars to
\(\hK\).

\subsection{The super Yetter--Drinfeld structure}
\label{sec:yd}

Introduce the group-like element
\begin{equation}\label{eq:Un}
 U_n=Q^{-2n\beta H_1+2n(\beta+1)H_2}.
\end{equation}
For homogeneous \(y\in\mathcal Y_n\) of bidegree \((a,b)\), define
\begin{equation}\label{eq:Yn-action}
\begin{aligned}
 y\tri H_1&=-by,& y\tri H_2&=-ay,\\
 y\tri x_1&=yx_1-Q^{-2n(\beta+1)}\chi(y,x_1)x_1y,\\
 y\tri x_2&=yx_2-Q^{2n\beta}\chi(y,x_2)x_2y.
\end{aligned}
\end{equation}
The boundary identities
\[
 v_n\tri x_1=v_n\tri x_2=0,\qquad
 ((x_1x_2)^{n-1}x_1)\tri x_2=v_n,\qquad
 ((x_2x_1)^{n-1}x_2)\tri x_1=t^{-1}v_n
\]
show directly that the action preserves \(\mathcal Y_n\).

For \(y\) of bidegree \((a,b)\), write the braided coproduct in
\(B_{\hK}\) as
\[
 \Delta_B(y)=\sum y^{(1)}\ot y^{(2)},\qquad
 \text{with bihomogeneous summands}.
\]
Then
\begin{equation}\label{eq:Yn-coaction}
 \delta_n^{\mathrm s}(y)
 =
 \sum (-1)^{|y^{(1)}||y^{(2)}|}
 \chi(y^{(2)},y^{(1)})^{-1}y^{(1)}\ot
 U_nQ^{-aH_1-bH_2}y^{(2)}.
\end{equation}
The explicit parity sign in \eqref{eq:Yn-coaction} cancels the Koszul
sign in \(\chi(y^{(2)},y^{(1)})^{-1}\); the super tensor category
already supplies the Koszul interchange.

\begin{proposition}\label{prop:YD}
Equations \eqref{eq:Yn-action} and \eqref{eq:Yn-coaction} make
\(\mathcal Y_n\) a right--right super Yetter--Drinfeld module over
\(\mathcal H^+\).
\end{proposition}

\begin{proof}
For \(y\) of bidegree \((a,b)\), set
\[
 \iota_n(y)=U_nQ^{-aH_1-bH_2}y\in\mathcal H^+.
\]
The action \eqref{eq:Yn-action} is the pullback along \(\iota_n\) of the
right super-adjoint action
\[
 A\tri_{\mathrm{ad}}h
 =\sum(-1)^{|A||h_{(1)}|}S(h_{(1)})Ah_{(2)}.
\]
Moreover,
\[
 \Delta(\iota_n(y))=(\iota_n\ot\id)\delta_n^{\mathrm s}(y).
\]
Thus \(\iota_n(\mathcal Y_n)\) is stable under the right adjoint action
and is a right coideal.  The standard adjoint/coideal calculation, with
the Koszul rule \eqref{eq:koszul}, gives the right--right super
Yetter--Drinfeld identity.  Counitality and coassociativity follow by
restricting those of \(\Delta\).
\end{proof}

The identification of the original operator \(T_n\) with the braiding
transported from this module is established in \cref{sec:rigid}; the
parameter and action conventions are checked in \cref{app:GK-conventions}.

\section{The completed paired double}\label{sec:double}

\subsection{The negative half and the Hopf pairing}

Define
\(\mathcal H^-\) using even generators \(\widehat H_1,\widehat H_2\),
and odd generators \(F_1,F_2\), with
\[
 [\widehat H_1,\widehat H_2]=0,\qquad
 [\widehat H_i,F_j]=a_{ji}F_j,\qquad F_i^2=0,
\]
and
\begin{equation}\label{eq:negative-Hopf}
\begin{aligned}
 \Delta(\widehat H_i)&=\widehat H_i\ot1+1\ot\widehat H_i,&
 S(\widehat H_i)&=-\widehat H_i,\\
 \Delta(F_i)&=F_i\ot Q^{\widehat H_i}+1\ot F_i,&
 S(F_i)&=-F_iQ^{-\widehat H_i}.
\end{aligned}
\end{equation}
The even Hopf pairing
\(\eta:\mathcal H^-\hot\mathcal H^+\to\hK\) is determined by
\begin{equation}\label{eq:eta-generators}
 \eta(\widehat H_i,H_j)=-\frac{a_{ji}}{\hbar},
 \qquad
 \eta(F_i,x_j)=\delta_{ij},
\end{equation}
all mixed Cartan--root pairings being zero, and by the right-dual order
\begin{equation}\label{eq:eta-extension}
\eta(fg,h)=\sum\eta(g,h_{(1)})\eta(f,h_{(2)}),\qquad
\eta(f,hk)=\sum\eta(f_{(1)},h)\eta(f_{(2)},k),
\end{equation}
with the Koszul sign understood in the second formula.

\subsection{Alternating-word bases, the Hopf pairing, and completion}\label{sec:completion}

For \(m\geq1\), let \(X_{i,m}\) be the alternating word of length \(m\)
in \(x_1,x_2\) beginning with \(x_i\), and let \(F_{i,m}\) be its
negative analogue.  The adjacent-square relations generate monomial
ideals, so the words that avoid adjacent squares form bases of the root
algebras.  These are exactly the alternating words.  Their degree pieces
are
\[
 B_m^+=\hK X_{1,m}\oplus\hK X_{2,m},\qquad
 B_m^-=\hK F_{1,m}\oplus\hK F_{2,m}.
\]
Set \(B_0^+=B_0^-=\hK1\), and let \(\mathcal C^+\) and
\(\mathcal C^-\) be the Cartan subalgebras generated by
\(H_1,H_2\) and \(\widehat H_1,\widehat H_2\), respectively:
their polynomial \(R\)-forms are \(\hbar\)-adically completed before
extension to \(\hK\), as specified in \cref{sec:conventions}.
The commuting Cartan generators act diagonally on the alternating
words.  The corresponding commuting weight derivations preserve the
adjacent-square ideals, so each half is the semidirect product of its
root algebra with its Cartan algebra.  Thus placing the Cartan
coefficients on the left gives unique normal forms and topological
vector-space decompositions
\[
 \mathcal H^\pm
 \cong\widehat{\bigoplus}_{m\geq0}\mathcal C^\pm B_m^\pm.
\]

\begin{lemma}\label{lem:pairing-descends}
The formulas \eqref{eq:eta-generators}--\eqref{eq:eta-extension} descend
uniquely to a continuous even Hopf pairing on
\(\mathcal H^-\hot\mathcal H^+\), and
\(\eta(B_m^-,B_l^+)=0\) for \(m\neq l\).
\end{lemma}

\begin{proof}
In the free Hopf superalgebras, the odd--odd terms in the coproduct cancel
by the Koszul rule, so the square relations on either side annihilate the
relation ideal on the other.  Primitive Cartan elements preserve Cartan
commutativity, while
\(\eta(Q^{\widehat H_i},Q^{H_j})=Q^{-a_{ji}}\) preserves the fixed-weight
Cartan--root relations.  Hence both relation ideals lie in the radicals
of the free pairing, and the pairing descends.  Root-degree orthogonality
follows from the skew-primitive coproducts.
\end{proof}

The Hopf pairing defines the right--right paired double
\[
 \mathfrak D_{\mathrm{rr},\eta}=\mathcal H^-\bowtie_\eta\mathcal H^+.
\]
The cross relations in the double are
\begin{equation}\label{eq:double-cross}
\begin{aligned}
 [H_i,\widehat H_j]&=0,&
 [H_i,F_j]&=-a_{ij}F_j,\\
 [\widehat H_i,x_j]&=-a_{ji}x_j,&
 x_iF_j+F_jx_i&=\delta_{ij}(Q^{H_i}-Q^{\widehat H_i}).
\end{aligned}
\end{equation}
Put \(P_r(Q)=\prod_{s=1}^r(1-Q^{-2s})\), with \(P_0(Q)=1\), and let
\[
 M_m(i,j)=\eta(F_{i,m},X_{j,m}).
\]

\begin{proposition}[Root-degree perfectness]\label{prop:perfect-pairing}
For \(r\geq0\) and \(r\geq1\), respectively,
\begin{equation}\label{eq:pairing-matrices}
\begin{aligned}
 M_{2r+1}&=(-1)^rP_r(Q)I_2,\\
 M_{2r}&=(-1)^rP_{r-1}(Q)
 \begin{pmatrix}1&-Q^{-r}\\-Q^{-r}&1\end{pmatrix}.
\end{aligned}
\end{equation}
Consequently,
\[
 \det M_{2r+1}=P_r(Q)^2,\qquad
 \det M_{2r}=P_{r-1}(Q)P_r(Q),
\]
and \(\eta|_{B_m^-\ot B_m^+}\) is perfect for every \(m\).
\end{proposition}

\begin{proof}
Apply \eqref{eq:eta-extension} recursively to the last letter of each
alternating word.  The two possible starting letters give a two-state
recurrence.  At odd length its off-diagonal entries cancel, while at even
length the two possible terminal pairings differ by \(Q^{-r}\).  Starting
from \(M_1=I_2\) and
\[
 M_2=\begin{pmatrix}-1&Q^{-1}\\Q^{-1}&-1\end{pmatrix}
\]
gives \eqref{eq:pairing-matrices} by induction.  Taking determinants and
using \(P_r=P_{r-1}(1-Q^{-2r})\) proves the claim.  Every \(P_r(Q)\) is
nonzero in \(\hK=\kk((\hbar))\).
\end{proof}

We use the directed root-degree completion
\begin{equation}\label{eq:directed-completion}
 \widehat{\mathcal T}_{+,-}
 =
 \prod_{m\geq0}
 (\mathcal C^+B_m^+)\hot_{\hK}(\mathcal C^-B_m^-).
\end{equation}
The product completes only the common root degree; in each factor,
the Cartan tensor product is formed from the completed \(R\)-forms
before extension to \(\hK\).  Root degrees add under multiplication, so
every fixed degree of a product contains only finitely many summands.

\begin{theorem}[Completed universal \(R\)-matrix]\label{thm:universal-R}
The signed right-dual coevaluation series of the paired double
\(\mathfrak D_{\mathrm{rr},\eta}\) belongs to
\eqref{eq:directed-completion} and is invertible there.  Define
\(\mathcal R_{\mathrm{rr}}\) to be its inverse.  Then
\begin{align}
 \mathcal R_{\mathrm{rr}}\Delta(d)
 &=\Delta^{\mathrm{op}}(d)\mathcal R_{\mathrm{rr}},
 \label{eq:R-intertwines}\\
 (\Delta\ot\id)(\mathcal R_{\mathrm{rr}})
 &=\mathcal R_{13}\mathcal R_{23},&
 (\id\ot\Delta)(\mathcal R_{\mathrm{rr}})
 &=\mathcal R_{13}\mathcal R_{12}.
\label{eq:R-hexagons}
\end{align}
\end{theorem}

\begin{proof}
\Cref{prop:perfect-pairing} supplies a dual basis in each root degree.
The super right-dual coevaluation contributes the additional sign
\((-1)^m\) in degree \(m\).  The resulting series belongs to
\eqref{eq:directed-completion}, and its inverse is defined recursively
degree by degree.  Dualizing multiplication and coproduct through
\eqref{eq:eta-extension} gives
\[
\begin{aligned}
 \Delta(d)\mathcal R_{\mathrm{rr}}^{-1}
 &=\mathcal R_{\mathrm{rr}}^{-1}\Delta^{\mathrm{op}}(d),\\
 (\Delta\ot\id)(\mathcal R_{\mathrm{rr}}^{-1})
 &=\mathcal R_{23}^{-1}\mathcal R_{13}^{-1},\\
 (\id\ot\Delta)(\mathcal R_{\mathrm{rr}}^{-1})
 &=\mathcal R_{12}^{-1}\mathcal R_{13}^{-1}.
\end{aligned}
\]
Inverting these identities proves \eqref{eq:R-intertwines} and
\eqref{eq:R-hexagons}.  No identification of either half with the full
continuous dual of the other is required.
\end{proof}

The double action induced from the Yetter--Drinfeld structure agrees,
through \eqref{eq:right-braiding}, with the canonical braiding
\(c_n^{\mathrm s}\).

\section{Root--center factorization}\label{sec:root-center}

\subsection{Root and central variables}

Define
\begin{equation}\label{eq:root-center-Cartan}
 \mathsf h_i=\frac{H_i-\widehat H_i}{2},\qquad
 \mathsf c_i=H_i+\widehat H_i.
\end{equation}
The \(\mathsf c_i\) are central.  Normalize the odd root vectors by
\begin{equation}\label{eq:root-generators}
 E_i=Q^{-\mathsf c_i/4}x_i,\qquad
 G_i=Q^{-\mathsf c_i/4}F_i.
\end{equation}
This separates the multiplication but not yet the coproduct.

\subsection{The Abelian twist}

Because \(A^{-1}=A\), define
\begin{equation}\label{eq:Frc}
\mathcal F_{\rc}
=
\exp\!\left[
\frac{\hbar}{4}
\sum_{i,j=1}^2(A^{-1})_{ji}
\bigl(\mathsf h_i\ot\mathsf c_j-\mathsf c_j\ot\mathsf h_i\bigr)
\right].
\end{equation}
All entries in the exponent are commuting even primitive elements.
Therefore \(\mathcal F_{\rc}\) is a normalized Drinfeld twist.
The twisted paired double has coproduct
\[
 \Delta_{\rc}(a)=\mathcal F_{\rc}\Delta(a)\mathcal F_{\rc}^{-1}.
\]
Its universal \(R\)-matrix is obtained by twisting the universal
\(R\)-matrix \(\mathcal R_{\mathrm{rr}}\) of \cref{thm:universal-R}:
\begin{equation}\label{eq:twisted-universal-R}
 \mathcal R_{\rc}
 =\mathcal F_{\rc,21}\mathcal R_{\mathrm{rr}}
  \mathcal F_{\rc}^{-1}.
\end{equation}
Thus \(\mathcal R_{\rc}\) is the universal \(R\)-matrix associated
with \(\Delta_{\rc}\), on the same underlying algebra.

For the convention \(\Delta_F=F\Delta F^{-1}\), the twisted antipode
is \(S_F=u_FS u_F^{-1}\), where
\(u_F=m(\id\ot S)(F)\).  In the present case the Cartan elements in
\eqref{eq:Frc} are commuting, even, and primitive, so
\[
 m(\id\ot S)(\mathcal F_{\rc})
 =\exp\!\left[
 \frac{\hbar}{4}\sum_{i,j=1}^2(A^{-1})_{ji}
 (-\mathsf h_i\mathsf c_j+\mathsf c_j\mathsf h_i)
 \right]=1.
\]
Consequently,
\begin{equation}\label{eq:twist-antipode-unchanged}
 S_{\rc}=S.
\end{equation}
Thus the twist leaves the dual-module actions unchanged, and the
canonical duality comparison for the twist equivalence is the identity.
The tensorator itself is generally nontrivial, but its contribution
contracts to the identity in evaluation and coevaluation.  This twist
therefore introduces no additional duality gauge.

\begin{theorem}[Root--center separation]\label{thm:root-center}
After twisting, the root and central Hopf structures separate:
\begin{equation}\label{eq:Hopf-factor}
 \mathfrak D_{\mathrm{rr},\eta}^{\mathcal F_{\rc}}
 \cong\mathcal D_{\rootpart}\hot\mathcal Z_{\cen}.
\end{equation}
Moreover,
\[
 \mathcal D_{\rootpart}
 \cong
 \hK\hot_R U_\hbar^{\Pi_{\ao}}(\mathfrak{sl}(2|1)).
\]
Here the standard positive and negative Chevalley generators correspond
to \(E_i\) and \((Q-Q^{-1})^{-1}G_i\), respectively, and the Cartan
generators correspond to \(\mathsf h_i\).
\end{theorem}

For clarity, the complete root presentation used in this theorem is
\begin{equation}\label{eq:all-odd-root-presentation}
\begin{gathered}
 [\mathsf h_i,\mathsf h_j]=0,
 \qquad [\mathsf h_i,E_j]=a_{ij}E_j,
 \qquad [\mathsf h_i,G_j]=-a_{ij}G_j,\\
 E_1^2=E_2^2=G_1^2=G_2^2=0,
 \qquad
 E_iG_j+G_jE_i
 =\delta_{ij}(Q^{\mathsf h_i}-Q^{-\mathsf h_i}).
\end{gathered}
\end{equation}
No additional quadratic relation is imposed between \(E_1E_2\) and
\(E_2E_1\), or between \(G_1G_2\) and \(G_2G_1\).  Every central
generator commutes with every generator in
\eqref{eq:all-odd-root-presentation}.

\begin{proof}
Conjugating the four skew-primitive terms by
\(\mathcal F_{\rc}\) gives
\[
\begin{aligned}
 \Delta_{\rc}(E_i)&=E_i\ot1+Q^{\mathsf h_i}\ot E_i,\\
 \Delta_{\rc}(G_i)&=G_i\ot Q^{-\mathsf h_i}+1\ot G_i.
\end{aligned}
\]
The \(\mathsf h_i,\mathsf c_i\) remain primitive, so their
exponentials remain group-like.  Thus no central generator occurs in the coproduct of
a root generator.  Conversely, the \(\mathsf c_i\) and their exponentials
generate a central commutative Hopf algebra.  The multiplication relations
from \eqref{eq:double-cross} become the all-odd
\(\mathfrak{sl}(2|1)\) relations, with both simple roots isotropic and
Cartan matrix \(A\).  To check that the algebra map is an
isomorphism, use the inverse substitutions
\[
 H_i=\mathsf h_i+\tfrac12\mathsf c_i,\qquad
 \widehat H_i=-\mathsf h_i+\tfrac12\mathsf c_i,\qquad
 x_i=Q^{\mathsf c_i/4}E_i,\qquad
 F_i=Q^{\mathsf c_i/4}G_i.
\]
Centrality of the \(\mathsf c_i\) and
\eqref{eq:all-odd-root-presentation} recover all Cartan--root
commutators and square relations.  They also give
\[
 x_iF_j+F_jx_i
 =\delta_{ij}Q^{\mathsf c_i/2}
 (Q^{\mathsf h_i}-Q^{-\mathsf h_i})
 =\delta_{ij}(Q^{H_i}-Q^{\widehat H_i}),
\]
recovering \eqref{eq:double-cross}.  Hence the inverse substitutions
define an algebra map inverse to the generator change
\eqref{eq:root-center-Cartan}--\eqref{eq:root-generators}.
The coproduct formulas above make this a Hopf isomorphism, proving
\eqref{eq:Hopf-factor}.
\end{proof}

\subsection[Factorization of the universal R-matrix]
{Factorization of the universal \(R\)-matrix}

We now describe how \(\mathcal R_{\rc}\) separates under the Hopf
isomorphism \eqref{eq:Hopf-factor}.
Let \(\mathcal R_{\rootpart}^{\rc}\) denote the inverse of the signed
canonical root tensor after replacing alternating \(x\)- and \(F\)-words
by the corresponding \(E\)- and \(G\)-words.  This is the root-vector
factor; the complete universal \(R\)-matrix of
\(\mathcal D_{\rootpart}\) also includes its Cartan factor.
Define the root and central Cartan factors by
\[
\begin{aligned}
 \mathcal R_{\mathrm{Cart}}^{\rootpart}
 &=
 \exp\!\left[-\hbar(
 \mathsf h_1\ot\mathsf h_2+\mathsf h_2\ot\mathsf h_1)\right],\\
 \mathcal R_{\mathrm{Cart}}^{\cen}
 &=
 \exp\!\left[\frac{\hbar}{4}(
 \mathsf c_1\ot\mathsf c_2+\mathsf c_2\ot\mathsf c_1)\right].
\end{aligned}
\]

\begin{proposition}\label{prop:R-factor}
In the directed completion,
\begin{equation}\label{eq:R-factor}
 \mathcal R_{\rc}
 =
 \mathcal R_{\rootpart}^{\rc}
 \mathcal R_{\mathrm{Cart}}^{\rootpart}
 \mathcal R_{\mathrm{Cart}}^{\cen}.
\end{equation}
Under \eqref{eq:Hopf-factor}, the product
\(\mathcal R_{\rootpart}^{\rc}\mathcal R_{\mathrm{Cart}}^{\rootpart}\)
belongs to the completed tensor square of \(\mathcal D_{\rootpart}\)
and is its universal \(R\)-matrix, while
\(\mathcal R_{\mathrm{Cart}}^{\cen}\) belongs to the completed tensor
square of \(\mathcal Z_{\cen}\) and is its universal \(R\)-matrix.
Their embeddings into the completed tensor square of the twisted double
commute, and their product is \(\mathcal R_{\rc}\).
Thus \eqref{eq:Hopf-factor} separates the quasitriangular structures as
well as the Hopf algebra structures.
\end{proposition}

\begin{proof}
Substitution of
\(H_i=\mathsf h_i+\frac12\mathsf c_i\) and
\(\widehat H_i=-\mathsf h_i+\frac12\mathsf c_i\) into the untwisted
Cartan exponential yields pure-root, mixed, and pure-central terms.
Since the exponent of \(\mathcal F_{\rc}\) is antisymmetric,
\(\mathcal F_{\rc,21}=\mathcal F_{\rc}^{-1}\); the twist factors in
\eqref{eq:twisted-universal-R} cancel the mixed terms exactly.
They simultaneously turn every root tensor
\(X_{i,m}\ot F_{j,m}\) into \(E_{i,m}\ot G_{j,m}\).
This proves \eqref{eq:R-factor}.  The universal \(R\)-identities follow
by twisting \cref{thm:universal-R}.
Under \eqref{eq:Hopf-factor}, apply the Hopf projections
\(\id\ot\varepsilon_{\cen}\) and
\(\varepsilon_{\rootpart}\ot\id\) to both tensor legs of these
identities.  The first projection removes the central Cartan factor,
leaving \(\mathcal R_{\rootpart}^{\rc}
\mathcal R_{\mathrm{Cart}}^{\rootpart}\); the second removes the root
factors, leaving \(\mathcal R_{\mathrm{Cart}}^{\cen}\).
Hence each is a universal \(R\)-matrix for the corresponding Hopf
factor.  Centrality gives the asserted commutation.
\end{proof}

Consequently, on modules that split into root and central factors,
the braiding splits into the corresponding two contributions.

\section{Identification of the coloring module}\label{sec:module}

\subsection{The root and central factors}

On a homogeneous \(y\in\mathcal Y_n\) of degree \((a,b)\), the root and
central Cartan elements act by
\begin{equation}\label{eq:root-center-actions}
\begin{aligned}
 y\tri\mathsf h_1&=(n(\beta+1)-b)y,&
 y\tri\mathsf h_2&=-(n\beta+a)y,\\
 y\tri\mathsf c_1&=-2n(\beta+1)y,&
 y\tri\mathsf c_2&=2n\beta y.
\end{aligned}
\end{equation}
Let \(W_{n,\beta}^{\ao}\) be the simple all-odd right highest-weight module with
even highest vector \(w_{n,\beta}\) satisfying
\begin{equation}\label{eq:highest-weight}
 w_{n,\beta}\tri E_1=w_{n,\beta}\tri E_2=0,\qquad
 w_{n,\beta}\tri\mathsf h_1=n\beta w_{n,\beta},\qquad
 w_{n,\beta}\tri\mathsf h_2=-n(\beta+1)w_{n,\beta}.
\end{equation}
Let \(L_{n,\beta}\) be the one-dimensional central module with characters
\[
 \mathsf c_1\mapsto-2n(\beta+1),\quad
 \mathsf c_2\mapsto2n\beta.
\]

\subsection{Module identification and the braiding factors}

\begin{theorem}[Module identification]\label{thm:object}
For \(\beta\neq0,-1\),
\begin{equation}\label{eq:object}
 \mathcal Y_n
 \cong W_{n,\beta}^{\ao}\boxtimes L_{n,\beta}
\end{equation}
as a right module over the twisted paired double.  The root factor
\(W_{n,\beta}^{\ao}\) is simple and has dimension \(4n\).
\end{theorem}

\begin{proof}
By \eqref{eq:Yn-action} and \eqref{eq:root-center-actions}, the nonzero
even vector \(v_n\) is killed by \(E_1,E_2\) and has the highest weight
\eqref{eq:highest-weight}.  Let
\(M=v_n\tri\mathcal D_{\rootpart}\subseteq\mathcal Y_n\).
As a cyclic highest-weight module, \(M\) has the simple highest-weight
quotient \(W_{n,\beta}^{\ao}\).  The dimension formula
\eqref{eq:typical-root-dimension}, established independently in
\cref{sec:module-typicality} below, and the basis
\eqref{eq:Yn-basis} give
\[
 4n=\dim_{\hK}W_{n,\beta}^{\ao}
 \leq\dim_{\hK}M
 \leq\dim_{\hK}\mathcal Y_n=4n.
\]
Thus \(M=\mathcal Y_n\) and the quotient map is an isomorphism of
root modules.  Finally, the central characters in
\eqref{eq:root-center-actions}, together with
\cref{thm:root-center}, supply the factor \(L_{n,\beta}\), proving
\eqref{eq:object}.
\end{proof}

Fixing a nonzero vector in \(L_{n,\beta}\), the identification
\eqref{eq:object} induces a root-module isomorphism
\[
 \Phi_n:\mathcal Y_n\xrightarrow{\sim}W_{n,\beta}^{\ao}.
\]

The factorization \eqref{eq:R-factor} now gives two contributions to
the braiding.  Define the root braiding by
\[
 R'_{n,\beta}
 =c^{\ao}_{W_{n,\beta}^{\ao},\,W_{n,\beta}^{\ao}}.
\]
By \eqref{eq:right-braiding}, it is obtained from the inverse of the
root universal \(R\)-matrix.  The inverse of the central Cartan factor
acts on \(L_{n,\beta}\ot L_{n,\beta}\), and hence on the central
factors of \(\mathcal Y_n\ot\mathcal Y_n\), by the scalar
\begin{equation}\label{eq:lambda}
 \lambda_{n,\beta}^{\cen}=Q^{2n^2\beta(\beta+1)}.
\end{equation}
Indeed, substituting the characters
\(\mathsf c_1=-2n(\beta+1)\) and \(\mathsf c_2=2n\beta\)
into \((\mathcal R_{\mathrm{Cart}}^{\cen})^{-1}\) gives this value.

\subsection{Odd reflection and typicality}
\label{sec:module-typicality}

We now determine the dimension of the abstract simple module
\(W_{n,\beta}^{\ao}\), independently of \(\mathcal Y_n\).
Set \(\alpha=n\beta\) and reflect the all-odd simple system at the
first isotropic root.  The distinguished Cartan data are
\[
\gamma_1=\alpha_1+\alpha_2,\qquad
\gamma_2=-\alpha_1,\qquad
h_1^{\mathrm{dist}}=\mathsf h_1+\mathsf h_2,\qquad
h_2^{\mathrm{dist}}=-\mathsf h_1.
\]
Write \(w=w_{n,\beta}\).  The relation
\eqref{eq:all-odd-root-presentation} gives
\[
 (w\tri G_1)\tri E_1=(Q^{n\beta}-Q^{-n\beta})w.
\]
Thus \(w^{\mathrm{dist}}=w\tri G_1\neq0\) when \(\beta\neq0\).
The odd-reflection calculation
\cite{Serganova,Yamane} makes this an odd highest vector for the
reflected Borel.  After the right-to-left conversion, its weights are
\[
 h_1^{\mathrm{dist}}\cdot w^{\mathrm{dist}}=(n-1)w^{\mathrm{dist}},
 \qquad
 h_2^{\mathrm{dist}}\cdot w^{\mathrm{dist}}=n\beta
 w^{\mathrm{dist}}.
\]
Here the conversion is \(a\cdot u=(-1)^{|a||u|}u\tri S(a)\) for
homogeneous \(a,u\); in particular it negates the Cartan weights.
If \(V_0(m,\alpha)\) denotes the simple distinguished highest-weight
module with an even highest vector, highest-weight classification gives
\begin{equation}\label{eq:odd-reflection-object}
 W_{n,\beta}^{\ao}\longleftrightarrow
 \Pi V_0(n-1,n\beta).
\end{equation}

For \(m\in\mathbb Z_{\geq0}\), the distinguished module
\(V_0(m,\alpha)\) is typical when
\(\alpha(\alpha+m+1)\neq0\), and then has dimension \(4(m+1)\).
These facts follow from the classical typicality criterion and character
formula in \cite[\S1.1, equation~(2)]{GeerPatureau}, together with the
character-preserving formal quantization recalled in
\cite[\S1.2]{GeerPatureau}; see also
\cite[\S3.3, equation~(39), and Lemma~3.4]{ColoredLG} for the
quantum-module formulation.
We use these formulas after extension to \(\hK\), with \(Q=e^\hbar\),
so there is no root-of-unity specialization.
For \(m=n-1\) and \(\alpha=n\beta\), the typicality condition is
\[
 \alpha(\alpha+m+1)=n^2\beta(\beta+1)\neq0.
\]
Consequently \(W_{n,\beta}^{\ao}\) is finite dimensional and
\begin{equation}\label{eq:typical-root-dimension}
 \dim_{\hK}W_{n,\beta}^{\ao}=4n
 \qquad(\beta\neq0,-1).
\end{equation}
The parity shift in \eqref{eq:odd-reflection-object} records our even
all-odd highest vector; the associated tangle conventions are compared
in \cref{sec:LG}.

\section{Oriented crossings and duality}\label{sec:rigid}

\subsection{Integer-colored braid gauges and transported duality}
\label{sec:oriented-fixed-color}

We first identify the positive crossing \(T_n\) with a gauge transform
of the canonical super Yetter--Drinfeld braiding.  The same gauge extends
to dual strands and cups and caps; \cref{sec:standard-duality} compares
the transported maps with standard partial transposes.

Fix \(\beta\in\kk\setminus\{0,-1\}\).  We construct a single gauge
for braid generators and oriented rigid transport with arbitrary integer
colors.  Positive colors label the modules constructed above, negative
colors label their left categorical duals, and zero labels the tensor unit.

For \(n\geq1\), write
\[
 \mathcal Y_{-n}:=\mathcal Y_n^\vee
 =\operatorname{Hom}_{\hK}(\mathcal Y_n,\hK).
\]
The negative index records categorical duality, rather than extending
the preceding construction to negative \(n\).  Also set
\(\mathcal Y_0:=\hK\), concentrated in even parity and bidegree
\((0,0)\), with the counit action.  Thus \(\mathcal Y_a\) is defined
for every \(a\in\mathbb Z\), with
\(\dim_{\hK}\mathcal Y_a=4|a|\) for \(a\ne0\) and
\(\dim_{\hK}\mathcal Y_0=1\).
For homogeneous \(f\in\mathcal Y_{-n}\), \(h\) in the paired double,
and \(y\in\mathcal Y_n\), its right action is
\begin{equation}\label{eq:dual-right-action}
 \langle f\tri h,y\rangle
 =
 (-1)^{|h|(|f|+1)}
 \langle f,y\tri S^{-1}(h)\rangle.
\end{equation}
This is the unique action for which evaluation is a right-module
morphism.  If \(\{e_a\}\) is a homogeneous basis and \(\{e^a\}\) its
homogeneous dual basis, we use
\begin{equation}\label{eq:dual-degree-cups}
 \deg e^a=-\deg e_a,\qquad
 \ev_n(f\ot y)=f(y),\qquad
 \coev_n(1)=\sum_a e_a\ot e^a.
\end{equation}
The sign in \eqref{eq:dual-right-action} and the use of \(S^{-1}\) are
forced by the left-dual convention.  With them, \(\ev_n\) and
\(\coev_n\) are module morphisms and satisfy the two super snake
identities.

Taking the left categorical dual of the module identification
\eqref{eq:object} gives, for \(n>0\),
\begin{equation}\label{eq:negative-object-factorization}
 \mathcal Y_{-n}
 \cong (W_{n,\beta}^{\ao})^\vee
       \boxtimes (L_{n,\beta})^\vee.
\end{equation}
This is an identification of modules for the twisted paired double.
By \eqref{eq:twist-antipode-unchanged}, its left dual agrees with the
one defined in \eqref{eq:dual-right-action}; no additional twist
duality gauge is required.  The root factor is the \(4n\)-dimensional
simple dual module, while the central factor is one-dimensional with
characters
\[
 \mathsf c_1\mapsto 2n(\beta+1),\qquad
 \mathsf c_2\mapsto-2n\beta.
\]
Thus negative colors are obtained by duality, without substituting a
negative integer into the highest-weight construction of \cref{sec:module}.

For every integer color \(a\), set
\begin{equation}\label{eq:oriented-objects-charges}
 U_a=Q^{-2a\beta H_1+2a(\beta+1)H_2}.
\end{equation}
Then \(U_{a+b}=U_aU_b\), \(U_{-a}=U_a^{-1}\), and \(U_0=1\).
For \(n>0\), the inverse charge on \(\mathcal Y_{-n}\) agrees with
the antipode in \eqref{eq:dual-right-action}.
If \(x\in\mathcal Y_a\) has bidegree \(p=(p_1,p_2)\), then
\[
 x\tri H_1=-p_2x,\qquad x\tri H_2=-p_1x,\qquad
 x\tri U_b=Q^{-2b(\beta+1)p_1+2b\beta p_2}x.
\]
These formulas hold for positive colors by \eqref{eq:Yn-action},
for negative colors by duality and the negative-degree convention,
and for zero by the counit action.

Let
\[
 \boldsymbol\alpha=(\alpha_1,\ldots,\alpha_r),
 \qquad
 \alpha_i\in\mathbb Z,\qquad
 \mathcal Y_{\boldsymbol\alpha}
 =\mathcal Y_{\alpha_1}\ot\cdots\ot \mathcal Y_{\alpha_r}.
\]
Extend \(\chi\) to nonzero bihomogeneous elements of
\(\bigoplus_{a\in\mathbb Z}\mathcal Y_a\) by the same formula
\eqref{eq:chi}, using the negative degrees on the dual from
\eqref{eq:dual-degree-cups}.  This covers pairs from any two colors;
parity is still the sum of the two degree coordinates modulo two.  For bihomogeneous
\(x_i\in \mathcal Y_{\alpha_i}\), define the oriented staircase gauge
\begin{equation}\label{eq:oriented-staircase}
\begin{aligned}
G_{\boldsymbol\alpha}^{\mathrm{or}}
(x_1\ot\cdots\ot x_r)
={}&
\left(\prod_{1\leq i<j\leq r}(-1)^{|x_i||x_j|}\chi(x_j,x_i)\right)\\
&{}\cdot
\bigotimes_{i=1}^r
\left(x_i\tri\prod_{j=i+1}^rU_{\alpha_j}\right).
\end{aligned}
\end{equation}
All charges lie in the same commutative torus, so the products in
\eqref{eq:oriented-staircase} are unambiguous.  Each gauge is diagonal
with nonzero powers of \(Q\) as eigenvalues, hence invertible.
A zero-colored tensor factor contributes no pairwise coefficient or
charge, so the gauge is compatible with the tensor-unit identifications.

For a right--right super Yetter--Drinfeld module, the checked braiding is
\begin{equation}\label{eq:YD-braiding}
 c_n^{\mathrm s}(a\ot b)
 =\sum(-1)^{|a||b_{(0)}|}
 b_{(0)}\ot(a\tri b_{(1)}).
\end{equation}
For \(y\in\mathcal Y_n\) of bidegree \((a,b)\), define
\[
 \varphi_n(y)=Q^{-2n(\beta+1)a+2n\beta b}y=y\tri U_n.
\]
The equality follows from \eqref{eq:Yn-action} and \eqref{eq:Un}.
For the all-positive word \((n,\ldots,n)\), the action on the
\(i\)-th factor in \eqref{eq:oriented-staircase} is
\(y_i\tri U_n^{r-i}=\varphi_n^{r-i}(y_i)\), with
\(\varphi_n^0=\id\).

For \(\alpha,\gamma\in\mathbb Z\), let
\[
 c_{\alpha,\gamma}^{\mathrm s}:
 \mathcal Y_\alpha\ot \mathcal Y_\gamma\longrightarrow \mathcal Y_\gamma\ot \mathcal Y_\alpha
\]
be the canonical braiding of the right-module category; in particular,
\(c_{n,n}^{\mathrm s}=c_n^{\mathrm s}\) for \(n>0\).
Define the colored \(T\)-crossings by
\begin{equation}\label{eq:oriented-T-crossing}
 T_{\alpha,\gamma}
 =
 G_{(\gamma,\alpha)}^{\mathrm{or}}\,
 c_{\alpha,\gamma}^{\mathrm s}\,
 (G_{(\alpha,\gamma)}^{\mathrm{or}})^{-1}.
\end{equation}
For each \(n>0\), \eqref{eq:oriented-T-crossing} includes
\(T_{n,n}\), \(T_{-n,n}\),
\(T_{n,-n}\), and \(T_{-n,-n}\).
The identification \(T_{n,n}=T_n\) with the original
Garoufalidis--Kashaev operator is verified below.

\begin{proposition}[Oriented braid-groupoid transport]
\label{prop:oriented-transport}
For \(r\geq2\) and \(1\leq i<r\), let
\(s_i\boldsymbol\alpha\) be obtained by exchanging the \(i\)-th and
\((i+1)\)-st entries.  Then
\begin{equation}\label{eq:oriented-local-transport}
 (T_{\boldsymbol\alpha})_i
 =
 G_{s_i\boldsymbol\alpha}^{\mathrm{or}}\,
 (c_{\boldsymbol\alpha}^{\mathrm s})_i\,
 (G_{\boldsymbol\alpha}^{\mathrm{or}})^{-1}.
\end{equation}
Under the identification of \cref{app:GK-conventions},
\(T_{n,n}=T_n\) for every \(n>0\).  Consequently, the operators in
\eqref{eq:oriented-T-crossing} extend the Garoufalidis--Kashaev braid
operators and satisfy the colored oriented braid-groupoid relations.
Crossings involving color zero are the canonical tensor-unit maps.
\end{proposition}

\begin{proof}
We first identify the positive crossing with the original operator.

Let \(a,b\in\mathcal Y_n\) be bihomogeneous.  For the
\(B_{\hK}\)-action, we use the same symbol \(\tri\) for the
Garoufalidis--Kashaev twisted right-adjoint action
\(a\tri h=\lambda_R(a\ot h)\); see \cref{app:GK-action}.  Decompose
\(\Delta_B(b)=\sum b^{(1)}\ot b^{(2)}\) into bihomogeneous components.
Expanding \eqref{eq:YD-braiding} using \eqref{eq:Yn-coaction} gives
\[
 c_n^{\mathrm s}(a\ot b)
 =\sum(-1)^{|b^{(2)}|(|b^{(1)}|+|a|)}
 \chi(b^{(2)},b^{(1)})^{-1}
 \chi(b^{(1)},a)\chi(b^{(2)},a)\,
 b^{(1)}\ot(\varphi_n(a)\tri b^{(2)}).
\]
To compute the conjugation, first apply
\[
 (G_{(n,n)}^{\mathrm{or}})^{-1}(a\ot b)
 =(-1)^{|a||b|}\chi(b,a)^{-1}
 \varphi_n^{-1}(a)\ot b.
\]
This cancels \(\varphi_n\) in the right-adjoint term above.
The resulting second tensor factor has degree
\(\deg a+\deg b^{(2)}\).  Applying \(G_{(n,n)}^{\mathrm{or}}\) therefore
contributes the scalar
\[
 (-1)^{(|a|+|b^{(2)}|)|b^{(1)}|}
 \chi(a,b^{(1)})\chi(b^{(2)},b^{(1)})
\]
and applies \(\varphi_n\) to the first tensor factor.
Since \(|b|=|b^{(1)}|+|b^{(2)}|\), all parity signs cancel.
Also,
\(\chi(b,a)=\chi(b^{(1)},a)\chi(b^{(2)},a)\).
Thus
\begin{align*}
 &G_{(n,n)}^{\mathrm{or}}c_n^{\mathrm s}(G_{(n,n)}^{\mathrm{or}})^{-1}(a\ot b)\\
 &\quad=
 \sum\chi(a,b^{(1)})\,
 \varphi_n(b^{(1)})\ot(a\tri b^{(2)})\\
 &\quad=T_n(a\ot b),
\end{align*}
where the last equality is \eqref{eq:GK-T-Sweedler} under the
identification of \cref{app:GK-conventions}.

For an arbitrary oriented word, the pairwise coefficients involving
any other strand cancel because the adjacent braiding preserves the
total bidegree of its two inputs.  The charge factors cancel as well:
a charge on a later strand acts on both adjacent inputs through its
group-like coproduct and commutes through the module braiding, while
an earlier strand sees the unchanged product
\(U_{\alpha_i}U_{\alpha_{i+1}}\).
The remaining factors are exactly the two-strand gauges in
\eqref{eq:oriented-T-crossing}.  This proves
\eqref{eq:oriented-local-transport}, including the all-positive case.
Conjugating the canonical braid relations by the boundary gauges proves
the braid-groupoid relations.
\end{proof}

Set \(G_\varnothing^{\mathrm{or}}=\id_{\hK}\).
For each \(n>0\), use the dual pair \((\mathcal Y_{-n},\mathcal Y_n)\)
with evaluation and coevaluation from \eqref{eq:dual-degree-cups}.
Its transported cups and caps are
\begin{equation}\label{eq:transported-cups}
\begin{aligned}
 \ev_n^T&=\ev_n(G_{(-n,n)}^{\mathrm{or}})^{-1},\\
 \coev_n^T&=G_{(n,-n)}^{\mathrm{or}}\coev_n.
\end{aligned}
\end{equation}
The comparison with standard duality is given in the next subsection.

\subsection{Comparison with standard duality}
\label{sec:standard-duality}

The equality \(T_{n,n}=T_n\) identifies the positive crossing.
The mixed crossings, however, use the transported duality maps
\eqref{eq:transported-cups}.  We now compare them with the standard
partial transposes used in \cite[Section~2]{GK}.  Throughout this
comparison \(n>0\), and tensor products of the underlying linear maps
are ordinary tensor products.

For \(y\) of bidegree \((a,b)\), define
\[
 A_n(y)=Q^{2ab-2n(\beta+1)a+2n\beta b}y.
\]
Its ordinary dual \(A_n^*\) is even and invertible.  Direct substitution
in \eqref{eq:oriented-staircase} gives
\begin{equation}\label{eq:cup-gauge}
 \ev_n=\ev_n^T((A_n^*)^{-1}\ot\id),\qquad
 \coev_n=(\id\ot A_n^*)\coev_n^T.
\end{equation}
Thus \(A_n^*\) identifies the transported negative object with the
standard dual of \(\mathcal Y_n\).  The underlying vector space is the
same; it is the evaluation and coevaluation that differ.

To specify the transpose convention, put \(V=\mathcal Y_n\).
For a linear map \(f:V\ot V\to V\ot V\), define
\begin{equation}\label{eq:standard-partial-transpose}
 \widetilde f
 =(\ev_n\ot\id_V\ot\id_{V^*})
 (\id_{V^*}\ot f\ot\id_{V^*})
 (\id_{V^*}\ot\id_V\ot\coev_n).
\end{equation}
Its source is \(V^*\ot V\) and its target is \(V\ot V^*\).
In an ordinary dual basis this reads
\[
 \widetilde f(e^a\ot e_b)
 =\sum_{c,d}f_{bd}^{ac}\,e_c\ot e^d,
 \qquad f(e_b\ot e_d)=\sum_{a,c}f_{bd}^{ac}e_a\ot e_c.
\]
This is the partial transpose of \cite[Section~2.1]{GK}; no super flip
is inserted in this contraction.  In particular, the operation is
linear in \(f\).  Taking an inverse crossing and taking a partial
transpose are distinct operations.

For the double-negative crossings, bend a second strand, following
\cite[Section~2, equations~(10)--(11)]{KashaevLong}:
\[
 \widetilde{\widetilde f}
 =(\ev_n\ot\id_{V^*}\ot\id_{V^*})
 (\id_{V^*}\ot\widetilde f\ot\id_{V^*})
 (\id_{V^*}\ot\id_{V^*}\ot\coev_n)
 :V^*\ot V^*\longrightarrow V^*\ot V^*.
\]
Here the second contraction uses the source and target of
\(\widetilde f\) specified above.  Explicitly,
\[
 \widetilde{\widetilde f}(e^a\ot e^c)
 =\sum_{b,d}f_{db}^{ca}\,e^b\ot e^d.
\]
Thus this double partial transpose is the full dual with the tensor
factors reversed; no super flip is inserted.

\begin{lemma}[Comparison of all oriented crossings]
\label{lem:partial-transpose}
The positive-positive crossings are \(T_{n,n}^{\pm1}=T_n^{\pm1}\).
The remaining six crossings satisfy
\begin{align}
 \widetilde{T_n^{-1}}
 &=(\id_V\ot A_n^*)T_{-n,n}
 ((A_n^*)^{-1}\ot\id_V)
 ,\label{eq:standard-negative-positive}\\
 (\widetilde{T_n})^{-1}
 &=(A_n^*\ot\id_V)T_{n,-n}
 (\id_V\ot(A_n^*)^{-1})
 ,
 \label{eq:standard-positive-negative}\\
 (\widetilde{T_n^{-1}})^{-1}
 &=(A_n^*\ot\id_V)T_{-n,n}^{-1}
 (\id_V\ot(A_n^*)^{-1})
 ,
 \label{eq:standard-negative-positive-inverse}\\
 \widetilde{T_n}
 &=(\id_V\ot A_n^*)T_{n,-n}^{-1}
 ((A_n^*)^{-1}\ot\id_V)
 ,
 \label{eq:standard-positive-negative-inverse}\\
 \widetilde{\widetilde{T_n}}
 &=(A_n^*\ot A_n^*)T_{-n,-n}
 ((A_n^*)^{-1}\ot(A_n^*)^{-1})
 ,
 \label{eq:standard-negative-negative}\\
 \widetilde{\widetilde{T_n^{-1}}}
 &=(A_n^*\ot A_n^*)T_{-n,-n}^{-1}
 ((A_n^*)^{-1}\ot(A_n^*)^{-1})
 .
 \label{eq:standard-negative-negative-inverse}
\end{align}
In the first and fourth identities the source and target are
\(V^*\ot V\) and \(V\ot V^*\); in the second and third they are
\(V\ot V^*\) and \(V^*\ot V\).  The last two are endomorphisms of
\(V^*\ot V^*\).
Hence, after changing the negative-edge identification by \(A_n^*\),
all oriented crossings agree with the standard-duality operators
prescribed in \cite[Section~2]{GK} and
\cite[Section~2, equations~(16)--(19)]{KashaevLong}.
\end{lemma}

\begin{proof}
First, the transported mate identities may be read using ordinary
local tensor products for the cup and cap contractions at issue here.
Indeed, a cup inserts degrees \(p,-p\) and charges \(U_n,U_{-n}\).
In the pairwise gauge, its coefficients with every external strand
cancel, since \(\chi\) is multiplicative in the degrees and
\(U_nU_{-n}=1\).  Only the internal two-strand factor remains,
which is exactly \(\coev_n^T\); the same argument applies to a cap.
Consequently, tensoring a local cup or cap with identity maps on
external factors is compatible with conjugation by the full boundary
gauges.  Together with \eqref{eq:oriented-local-transport}, this shows
that each local crossing, cup, or cap is the corresponding canonical
map conjugated by its input and output boundary gauges.
These gauges cancel on composition, so the canonical snake and
braiding evaluation identities give the transported identities using
ordinary tensor products.

The braiding evaluation identity expresses \(T_{-n,n}\) as the
contraction \eqref{eq:standard-partial-transpose} of \(T_n^{-1}\),
with \(\ev_n,\coev_n\) replaced by \(\ev_n^T,\coev_n^T\).
Substituting \eqref{eq:cup-gauge} into that contraction gives
\[
 T_{-n,n}
 =(\id_V\ot(A_n^*)^{-1})
 \widetilde{T_n^{-1}}(A_n^*\ot\id_V),
\]
which proves \eqref{eq:standard-negative-positive}.
The other evaluation identity expresses \((T_{n,-n})^{-1}\) as the
same transported contraction of \(T_n\).  Substitution gives
\[
 (T_{n,-n})^{-1}
 =(\id_V\ot(A_n^*)^{-1})\widetilde{T_n}
 (A_n^*\ot\id_V).
\]
Inverting proves \eqref{eq:standard-positive-negative} and also
establishes the required invertibility of both partial transposes.
Taking inverses gives
\eqref{eq:standard-negative-positive-inverse} and
\eqref{eq:standard-positive-negative-inverse}.

Bending both strands gives the categorical dual of the positive
crossing.  The dual of a tensor product reverses its factor order,
so, in the ordinary dual bases fixed above, this map is exactly
\(\widetilde{\widetilde{T_n}}\), as the displayed coefficient formula
shows.  Each negative boundary factor is identified with the
standard dual by \(A_n^*\), giving
\eqref{eq:standard-negative-negative}.  Taking its inverse gives
\eqref{eq:standard-negative-negative-inverse}, since the reversed
tensor dual commutes with inversion:
\[
 \widetilde{\widetilde{T_n^{-1}}}
 =\bigl(\widetilde{\widetilde{T_n}}\bigr)^{-1}.
\]
\end{proof}

For clarity, temporarily write \(J^{\mathrm{tr}}_{T_n}(D)\) for the
normal long-knot operator computed with the transported crossings and
cups, and \(J^{\mathrm{GK}}_{T_n}(D)\) for the standard-duality
state sum of \cite[Section~2.2]{GK}.  Change the identification on
every negative edge by \(A_n^*\), leaving positive edges unchanged.
The lemma and \eqref{eq:cup-gauge} convert all local oriented crossings,
cups, and caps to their standard-duality counterparts.
On every internal edge the identification and its inverse
cancel.  Both external edges are positive, so
\begin{equation}\label{eq:standard-state-sum-comparison}
 J^{\mathrm{tr}}_{T_n}(D)=J^{\mathrm{GK}}_{T_n}(D).
\end{equation}
We henceforth denote this common operator by \(J_{T_n}(D)\).
This is equality of the assembled operators, not equality of the
individual mixed-crossing matrices before changing the dual
identification.

For later use, the same bending argument applies to the root braiding
\(R'_{n,\beta}\): its standard partial transposes and those of its
inverse are invertible.  For any nonzero scalar \(\mu\), linearity and
inversion give
\[
 \widetilde{\mu R}=\mu\widetilde R,\qquad
 \widetilde{(\mu R)^{-1}}=\mu^{-1}\widetilde{R^{-1}},\qquad
 (\widetilde{\mu R})^{-1}=\mu^{-1}(\widetilde R)^{-1}.
\]
The inverse scalar in a mixed crossing therefore comes from taking
an inverse, not from partial transpose itself.

\section{Long-knot transport and the Links--Gould comparison}
\label{sec:LG}

\subsection{Compatibility with the root--center twist}
\label{sec:oriented-root-center}

The preceding gauge connects \(T_n\) with the untwisted braiding.  We now
include the tensorator of \(\mathcal F_{\rc}\) on the same oriented words.
If \(x\in \mathcal Y_\alpha\) has bidegree \(p=(p_1,p_2)\), then
\begin{equation}\label{eq:oriented-Cartan-weights}
\begin{aligned}
 x\tri\mathsf h_1
 &=(\alpha(\beta+1)-p_2)x,&
 x\tri\mathsf h_2
 &=(-\alpha\beta-p_1)x,\\
 x\tri\mathsf c_1
 &=-2\alpha(\beta+1)x,&
 x\tri\mathsf c_2
 &=2\alpha\beta x.
\end{aligned}
\end{equation}
For \(\alpha,\gamma\in\mathbb Z\), define
\begin{equation}\label{eq:oriented-F-exponent}
\begin{aligned}
\Gamma_{\alpha\mid\gamma}(p,q)
=\frac12\bigl[
&\gamma((\beta+1)p_1-\beta p_2)\\
&+\alpha(\beta q_2-(\beta+1)q_1)
\bigr].
\end{aligned}
\end{equation}
Let \(\mathcal F_{\rc}^{[r]}\) be the coherent \(r\)-fold tensorator of
the normalized twist and put
\begin{equation}\label{eq:oriented-F-tensorator}
 \mathsf F_{\boldsymbol\alpha}^{\mathrm{or}}
 =
 \rho_{\mathcal Y_{\boldsymbol\alpha}}(\mathcal F_{\rc}^{[r]}).
\end{equation}
Equations \eqref{eq:Frc} and \eqref{eq:oriented-Cartan-weights} give the
explicit diagonal action
\begin{equation}\label{eq:oriented-F-action}
\mathsf F_{\boldsymbol\alpha}^{\mathrm{or}}
(x_1\ot\cdots\ot x_r)
=
Q^{\sum_{i<j}
\Gamma_{\alpha_i\mid\alpha_j}(p_i,p_j)}
x_1\ot\cdots\ot x_r.
\end{equation}
The negative-color signs in this formula come from categorical duality,
not from an additional convention.

Let \((c_{\boldsymbol\alpha}^{\rc})_i\) be the braiding of the twisted
paired double.  Monoidal transport under a Drinfeld twist gives
\begin{equation}\label{eq:oriented-twisted-braiding}
 (c_{\boldsymbol\alpha}^{\rc})_i
 =
 (\mathsf F_{s_i\boldsymbol\alpha}^{\mathrm{or}})^{-1}
 (c_{\boldsymbol\alpha}^{\mathrm s})_i
 \mathsf F_{\boldsymbol\alpha}^{\mathrm{or}}.
\end{equation}
Define the combined gauge
\begin{equation}\label{eq:combined-oriented-gauge}
 \mathsf G_{\boldsymbol\alpha}^{\rc,\mathrm{or}}
 =
 G_{\boldsymbol\alpha}^{\mathrm{or}}
 \mathsf F_{\boldsymbol\alpha}^{\mathrm{or}}.
\end{equation}
Substitution in \eqref{eq:oriented-local-transport} yields
\begin{equation}\label{eq:oriented-combined-transport}
 (T_{\boldsymbol\alpha})_i
 =
 \mathsf G_{s_i\boldsymbol\alpha}^{\rc,\mathrm{or}}
 (c_{\boldsymbol\alpha}^{\rc})_i
 (\mathsf G_{\boldsymbol\alpha}^{\rc,\mathrm{or}})^{-1}.
\end{equation}
This formula relates the integer-colored \(T\)-crossings to the
braiding of the twisted paired double through the combined gauge.

For each \(n>0\), the twist leaves the standard evaluation and
coevaluation unchanged:
\begin{equation}\label{eq:root-center-cups}
 \ev_n\mathsf F_{(-n,n)}^{\mathrm{or}}=\ev_n,\qquad
 (\mathsf F_{(n,-n)}^{\mathrm{or}})^{-1}\coev_n=\coev_n.
\end{equation}
Indeed, their paired degrees give
\(\Gamma_{-n\mid n}(-p,p)=\Gamma_{n\mid -n}(p,-p)=0\),
in agreement with \eqref{eq:twist-antipode-unchanged}.
Using \eqref{eq:combined-oriented-gauge} and
\eqref{eq:transported-cups}, the combined gauge therefore satisfies
\begin{align}
 \ev_n(\mathsf G_{(-n,n)}^{\rc,\mathrm{or}})^{-1}
 &=\ev_n^T,\label{eq:root-center-evaluation-cancels}\\
 \mathsf G_{(n,-n)}^{\rc,\mathrm{or}}\coev_n
 &=\coev_n^T.\label{eq:root-center-coevaluation-cancels}
\end{align}
These identities provide the local crossing, cup, and cap comparisons
needed for long-knot transport.

\subsection{The normal long-knot operator}

Let \(D\) be a normal long-knot diagram and let \(J_R(D)\) be its
endomorphism-valued state sum using standard duality as in
\cite[Section~2]{GK}.  For \(R=T_n\), its agreement with the
transported construction is \eqref{eq:standard-state-sum-comparison}.
Composing the local comparison
at every crossing, cup, and cap gives the following global statement.

\begin{theorem}[Long-knot transport]\label{thm:rigid}
For every normal long-knot diagram \(D\),
\begin{equation}\label{eq:long-transport}
 \Phi_nJ_{T_n}(D)\Phi_n^{-1}
 =
 (\lambda_{n,\beta}^{\cen})^{\Writhe(D)}
 J_{R'_{n,\beta}}(D).
\end{equation}
In particular, the two operators are conjugate whenever \(D\) is
balanced, that is, \(\Writhe(D)=0\).
\end{theorem}

\begin{proof}
By \cref{lem:partial-transpose},
\eqref{eq:standard-state-sum-comparison}, and the combined local
transport identities above, the boundary gauges cancel when the
elementary pieces of \(D\) are composed.  The remaining external
identification is \(\Phi_n\).  Root--center factorization contributes
\(\lambda_{n,\beta}^{\cen}\) at each positive crossing and its inverse
at each negative crossing, by the scalar rules in
\cref{sec:standard-duality}.  Their product is
\((\lambda_{n,\beta}^{\cen})^{\Writhe(D)}\), proving the formula.
\end{proof}

\subsection{Balanced representatives}

For the knot comparison below, choose a balanced normal diagram \(D\)
representing \(\mathcal K\).  Such a representative is obtained by
composing a normal diagram with Kashaev's correction
\(\xi^{-\Writhe(D)/2}\): normal diagrams have even writhe and
\(\Writhe(\xi^k)=2k\)
\cite[Section~1, Remark~3, and Section~2, Theorem~1]{KashaevLong}.
The state sum on this corrected diagram is the normalized knot operator.
In the remainder of this section, \(D\) denotes a balanced normal
representative, so the central scalar in \cref{thm:rigid} is one.

\subsection{Scalarity}

The operator \(J_{R'_{n,\beta}}(D)\) is an even
\(U_\hbar^{\Pi_{\ao}}(\mathfrak{sl}(2|1))\)-module endomorphism of
\(W_{n,\beta}^{\ao}\): braidings, their inverses, cups, caps, tensor
products, and compositions are all module morphisms.

\begin{proposition}\label{prop:Schur}
\[
 \End_{U_\hbar^{\Pi_{\ao}}(\mathfrak{sl}(2|1))}
 (W_{n,\beta}^{\ao})
 =\hK\,\id.
\]
Consequently, there is a unique scalar
\(LG'_{n,\mathcal K}(Q^{n\beta},Q)\) such that
\[
 J_{R'_{n,\beta}}(D)
 =LG'_{n,\mathcal K}(Q^{n\beta},Q)\id.
\]
\end{proposition}

\begin{proof}
An endomorphism sends the one-dimensional highest-weight space to itself.
It therefore acts on \(w_{n,\beta}\) by a scalar.  Since
\(W_{n,\beta}^{\ao}\) is generated by \(w_{n,\beta}\), the same scalar
acts on the whole module.
\end{proof}

Combining \cref{thm:rigid,prop:Schur}, and recalling that
\(V_{n,\mathcal K}\) is the \((1,1)\)-entry of \(J_{T_n}(D)\), gives
\begin{equation}\label{eq:V-LGprime}
 V_{n,\mathcal K}(Q^{2n\beta+n},Q^2)
 =LG'_{n,\mathcal K}(Q^{n\beta},Q).
\end{equation}
In fact, the argument proves the stronger scalar identity
\[
 J_{T_n}(D)
 =LG'_{n,\mathcal K}(Q^{n\beta},Q)\id_{\mathcal Y_n}
\]
for balanced normal diagrams.  Thus scalarity is a theorem rather than
an input.

\subsection{Comparison with standard Links--Gould conventions}

It remains to compare the root-module scalar with the standard
Links--Gould normalization.

\begin{theorem}\label{thm:standard-comparison}
For every oriented knot \(\mathcal K\),
\begin{equation}\label{eq:LGprime-standard}
 LG'_{n,\mathcal K}(Q^{n\beta},Q)
 =
 LG_{\mathcal K}^{(n)}(Q^{-n\beta},Q^{-1})
 =
 LG_{\overline{\mathcal K}}^{(n)}(Q^{n\beta},Q).
\end{equation}
\end{theorem}

\begin{proof}
The odd reflection \eqref{eq:odd-reflection-object} identifies the
coloring object with \(\Pi V_0(n-1,n\beta)\).  Independence of the
Drinfeld--Jimbo quantization from the simple-root system, together with
the quantum-group/Kontsevich comparison, identifies the corresponding
normalized rigid-tangle scalars \cite{GeerCartan,GeerKontsevich}.
The parity shift contributes only a writhe factor, which is one on the
balanced representatives used here.

Standard Links--Gould conventions use left modules with braiding
\(\tau^{\mathrm s}(\pi\ot\pi)(\mathcal R)\), whereas
\eqref{eq:right-braiding} uses \(\mathcal R^{-1}\).  The right-to-left
conversion therefore reverses the crossings, producing the simultaneous
substitutions \((Q^{n\beta},Q)\mapsto(Q^{-n\beta},Q^{-1})\), or,
equivalently, taking the mirror with parameters fixed.  This convention
is uniform in \(n\) and agrees with
\cite[Theorem~1.5]{ColoredLG}, proving \eqref{eq:LGprime-standard}.
\end{proof}

\begin{proof}[Proof of \cref{thm:main}]
Substitute \eqref{eq:LGprime-standard} into \eqref{eq:V-LGprime}.
The scalar-identity assertion follows from \cref{prop:Schur}.
\end{proof}

\subsection{Concluding remarks}

The staircase gauge identifies the Nichols braid operator with the
super Yetter--Drinfeld braiding, while the root--center twist separates
the quantum-group contribution from the scalar
\(Q^{2n^2\beta(\beta+1)}\).  Balanced normalization removes this scalar.
The remaining long-knot operator acts on a simple typical module, so its
\((1,1)\)-entry determines the entire operator.

\appendix

\section*{Appendices}
\phantomsection
\addcontentsline{toc}{section}{Appendices}

\section{Comparison with the Garoufalidis--Kashaev conventions}
\label{app:GK-conventions}

We record the parameter and action conventions relating our construction
to \cite[Theorem~3.6 and Definition~7.5]{GK}.

\subsection{Parameters and the space \texorpdfstring{\(Y_n\)}{Yn}}
\label{app:GK-parameters}

To distinguish the notation of \cite{GK} from ours, add a superscript
\({\rm GK}\) to its parameters.  The required specialization is
\begin{equation}\label{eq:GK-dictionary}
 q^{\rm GK}=Q^2,\qquad
 q_{12}^{\rm GK}=q_{21}^{\rm GK}=-Q,\qquad
 t^{\rm GK}=Q^{2n\beta+n}=t.
\end{equation}
Indeed, the diagonal braiding determined by
\(q_{11}^{\rm GK}=q_{22}^{\rm GK}=-1\) and
\(q_{12}^{\rm GK}=q_{21}^{\rm GK}=-Q\) is exactly the braiding
\(\tau_\chi\) determined by \eqref{eq:chi}.  The parametrization following
\cite[Proposition~7.4]{GK} therefore gives
\begin{equation}\label{eq:GK-t1-t2}
 \begin{aligned}
 t_1^{\rm GK}
 &=\frac{1}{(q^{\rm GK})^{n/2}t^{\rm GK}}
   =Q^{-2n(\beta+1)},\\
 t_2^{\rm GK}
 &=\frac{t^{\rm GK}}{(q^{\rm GK})^{n/2}}
   =Q^{2n\beta}.
 \end{aligned}
\end{equation}
In particular, \(t_1^{\rm GK}t_2^{\rm GK}(q^{\rm GK})^n=1\).
The coefficient of the invariant vector in \cite[(110)--(113)]{GK} is
\begin{equation}\label{eq:GK-alpha}
 \alpha^{\rm GK}=t_2^{\rm GK}(-q_{12}^{\rm GK})^n
 =Q^{2n\beta+n}=t.
\end{equation}
Consequently,
\[
 v_{n,\alpha^{\rm GK}}=(x_1x_2)^n+\alpha^{\rm GK}(x_2x_1)^n=v_n,
\]
and the basis in \cite[Proposition~7.4]{GK} becomes exactly the basis
\eqref{eq:Yn-basis}.  Thus, after the specialization
\eqref{eq:GK-dictionary} and extension of scalars to \(\hK\), there is a
canonical identification
\begin{equation}\label{eq:GK-space-identification}
 Y_n^{\rm GK}\otimes_{\mathbb F}\hK\cong\mathcal Y_n
\end{equation}
inside \(B_{\hK}\).  Since \(Q=e^\hbar\), the parameter \(q^{\rm GK}=Q^2\)
is not a root of unity in the formal setting.  Moreover,
\(\beta\ne-1,0\) gives respectively
\(t_1^{\rm GK}\ne1\) and \(t_2^{\rm GK}\ne1\), so the hypotheses of
\cite[Proposition~7.4]{GK} hold.

\subsection{The right action and the original definition of
\texorpdfstring{\(T_n\)}{Tn}}
\label{app:GK-action}

The map \(\varphi_n\) used in \cref{sec:oriented-fixed-color} is the restriction of the
scaling automorphism of the braided Hopf algebra \(B_{\hK}\) determined by
\begin{equation}\label{eq:GK-scaling}
 \varphi_n(x_1)=Q^{-2n(\beta+1)}x_1,\qquad
 \varphi_n(x_2)=Q^{2n\beta}x_2.
\end{equation}
We also write
\[
 a\tri h=\lambda_R(a\ot h),\qquad a,h\in B_{\hK},
\]
for the twisted right-adjoint action of \cite[Theorem~3.5]{GK}.
Its restriction to \(\mathcal Y_n\) agrees with the action above
when the acting element belongs to \(B_{\hK}\); Cartan and paired-double
actions retain their previously specified domains.
On the generators it is
\begin{equation}\label{eq:GK-right-action}
\begin{aligned}
 a\tri x_1&=ax_1-Q^{-2n(\beta+1)}\chi(a,x_1)x_1a,\\
 a\tri x_2&=ax_2-Q^{2n\beta}\chi(a,x_2)x_2a,
\end{aligned}
\end{equation}
which is precisely the restriction used in \eqref{eq:Yn-action}.

Under \eqref{eq:GK-space-identification}, the operator called \(T_n\) in
\cite[Theorem~3.6 and Definition~7.5]{GK} is
\begin{equation}\label{eq:GK-T-definition}
 T_n^{\rm GK}
 =
 (\varphi_n\ot\lambda_R)
 (\tau_\chi\ot\id)
 (\id\ot\Delta_B)
 \big|_{\mathcal Y_n\ot\mathcal Y_n}.
\end{equation}
Thus, if \(a,b\) are bihomogeneous and
\[
 \Delta_B(b)=\sum b^{(1)}\ot b^{(2)}
 \quad\text{in bihomogeneous components},
\]
then
\begin{equation}\label{eq:GK-T-Sweedler}
 T_n^{\rm GK}(a\ot b)
 =\sum \chi(a,b^{(1)})\,
 \varphi_n(b^{(1)})\ot(a\tri b^{(2)}).
\end{equation}
This fixes both the tensor order and the right-action convention.

\subsection{Parameter dictionary}

\begin{center}
\begin{tabular}{@{}lll@{}}
\toprule
Symbol & Meaning & Relation\\
\midrule
\(Q\) & formal quantum parameter & \(Q=e^\hbar\)\\
\(\beta\) & highest-weight parameter & \(\beta\neq0,-1\)\\
\(\alpha\) & standard Links--Gould weight & \(\alpha=n\beta\)\\
\(t\) & first \(V_n\)-parameter & \(t=Q^{2n\beta+n}\)\\
\(q\) & second \(V_n\)-parameter & \(q=Q^2\)\\
\bottomrule
\end{tabular}
\end{center}

\section{First root degrees of the canonical tensor}

For reference, the inverses of the pairing matrices are
\[
\begin{aligned}
 M_{2r+1}^{-1}&=\frac{(-1)^r}{P_r(Q)}I_2,\\
 M_{2r}^{-1}&=\frac{(-1)^r}{P_r(Q)}
 \begin{pmatrix}1&Q^{-r}\\Q^{-r}&1\end{pmatrix}.
\end{aligned}
\]
The super right-dual convention multiplies degree \(m\) by \((-1)^m\).
Thus the signed root coevaluation begins with
\[
 1-x_1\ot F_1-x_2\ot F_2+\text{terms of root degree at least two}.
\]
Its inverse, the standard root factor of the universal \(R\)-matrix,
begins with
\[
 1+x_1\ot F_1+x_2\ot F_2+\text{terms of root degree at least two}.
\]
These signs are forced simultaneously by the right-dual evaluation and
the cross relation \(x_iF_i+F_ix_i=Q^{H_i}-Q^{\widehat H_i}\).

\section*{Acknowledgments}
\phantomsection
\addcontentsline{toc}{section}{Acknowledgments}

The author is deeply grateful to his advisor, Nicolai Reshetikhin, for
carefully reviewing this article, offering valuable suggestions, and providing
encouragement and support.  He also warmly thanks his family, friends, and
partner for standing by him during times of discouragement.

The author used \texttt{gpt5.6sol} to obtain the key ideas of the proof,
and used \texttt{gpt5.6sol} and \texttt{gpt6.1sol} in repeated rounds of
revision and polishing to make the exposition readable.  He also thanks
Hanwen Quan for advice on the use of AI.  The author takes full
responsibility for the content and for any remaining errors.

\begin{thebibliography}{99}
\raggedright

\bibitem{DeWitKauffmanLinks}
D.~De Wit, L.~H.~Kauffman, and J.~R.~Links,
\newblock On the Links--Gould invariant of links,
\newblock \emph{J. Knot Theory Ramifications} \textbf{8} (1999), no.~2,
165--199,
\newblock \href{https://doi.org/10.1142/S0218216599000110}
{doi:10.1142/S0218216599000110}.

\bibitem{ReshetikhinTuraev}
N.~Yu.~Reshetikhin and V.~G.~Turaev,
\newblock Ribbon graphs and their invariants derived from quantum groups,
\newblock \emph{Comm. Math. Phys.} \textbf{127} (1990), no.~1, 1--26,
\newblock \href{https://doi.org/10.1007/BF02096491}
{doi:10.1007/BF02096491}.

\bibitem{Kac}
V.~G.~Kac,
\newblock Lie superalgebras,
\newblock \emph{Adv. Math.} \textbf{26} (1977), no.~1, 8--96,
\newblock \href{https://doi.org/10.1016/0001-8708(77)90017-2}
{doi:10.1016/0001-8708(77)90017-2}.

\bibitem{Serganova}
V.~Serganova,
\newblock On generalizations of root systems,
\newblock \emph{Comm. Algebra} \textbf{24} (1996), no.~13, 4281--4299,
\newblock \href{https://doi.org/10.1080/00927879608825814}
{doi:10.1080/00927879608825814}.

\bibitem{Yamane}
H.~Yamane,
\newblock Quantized enveloping algebras associated with simple Lie
superalgebras and their universal \(R\)-matrices,
\newblock \emph{Publ. Res. Inst. Math. Sci.} \textbf{30} (1994), no.~1,
15--87,
\newblock \href{https://doi.org/10.2977/PRIMS/1195166275}
{doi:10.2977/PRIMS/1195166275}.

\bibitem{GK}
S.~Garoufalidis and R.~Kashaev,
\newblock Multivariable knot polynomials from braided Hopf algebras with
automorphisms,
\newblock \emph{Publ. Res. Inst. Math. Sci.} \textbf{62} (2026), 75--114,
\newblock \href{https://doi.org/10.4171/PRIMS/62-1-3}
{doi:10.4171/PRIMS/62-1-3}.

\bibitem{KashaevLong}
R.~Kashaev,
\newblock Invariants of long knots,
\newblock in \emph{Representation Theory, Mathematical Physics, and
Integrable Systems}, Progress in Mathematics \textbf{340},
Birkh\"auser, Cham, 2021, 431--451,
\newblock \href{https://arxiv.org/abs/1908.00118}{arXiv:1908.00118}.

\bibitem{ColoredLG}
S.~Garoufalidis, M.~Harper, R.~Kashaev, B.-M.~Kohli, and E.~Wagner,
\newblock On the colored Links--Gould polynomial,
\newblock arXiv:2509.10911, 2025,
\newblock \url{https://arxiv.org/abs/2509.10911}.

\bibitem{GeerCartan}
N.~Geer,
\newblock Some remarks on quantized Lie superalgebras of classical type,
\newblock \emph{J. Algebra} \textbf{314} (2007), no.~2, 565--580,
\newblock \href{https://doi.org/10.1016/j.jalgebra.2006.10.045}
{doi:10.1016/j.jalgebra.2006.10.045}.

\bibitem{GeerEK}
N.~Geer,
\newblock Etingof--Kazhdan quantization of Lie superbialgebras,
\newblock \emph{Adv. Math.} \textbf{207} (2006), no.~1, 1--38,
\newblock \href{https://doi.org/10.1016/j.aim.2005.11.005}
{doi:10.1016/j.aim.2005.11.005}.

\bibitem{GeerKontsevich}
N.~Geer,
\newblock The Kontsevich integral and quantized Lie superalgebras,
\newblock \emph{Algebr. Geom. Topol.} \textbf{5} (2005), 1111--1139,
\newblock \href{https://doi.org/10.2140/agt.2005.5.1111}
{doi:10.2140/agt.2005.5.1111}.

\bibitem{GeerPatureau}
N.~Geer and B.~Patureau--Mirand,
\newblock Multivariable link invariants arising from
\(\mathfrak{sl}(2|1)\) and the Alexander polynomial,
\newblock \emph{J. Pure Appl. Algebra} \textbf{210} (2007), no.~1,
283--298,
\newblock \url{https://arxiv.org/abs/math/0601291}.

\end{thebibliography}

\end{document}
